\section*{Chapitre \ref{ch:iv1:optimal-stopping} --- Arrêt optimal}
\markboth{\protect\omTocMark{}{Réponses~: Arrêt optimal}}{\protect\omTocMark{}{Réponses~: Arrêt optimal}}

\begin{solution}{iq:iv1:optimal-stopping:1}
Avec une relance, on garde 4, 5 ou 6, et le jeu vaut $4.25$. Avec deux relances, on garde d'abord 5 ou 6, et le jeu vaut $14/3 \approx 4.67$.

Raisonner à rebours\index{induction à rebours}. Le dernier lancer vaut $3.5$, donc avec une relance restante on garde tout lancer supérieur à $3.5$. La valeur vaut alors $\tfrac12 \cdot 5 + \tfrac12 \cdot 3.5 = 4.25$. Avec deux relances restantes, on compare le lancer à $4.25$, donc on ne garde que 5 et 6. La valeur est $\tfrac13 \cdot 5.5 + \tfrac23 \cdot 4.25 = 14/3$. Avec beaucoup de relances, la valeur grimpe vers 6.

\iqwrong{Utiliser le même seuil à chaque étape. Le seuil est la valeur de ce que l'on abandonne, et cette valeur est d'autant plus élevée qu'il reste davantage de relances.}
\end{solution}

\begin{solution}{iq:iv1:optimal-stopping:2}
Les valeurs sont $5/8$, $89/128 \approx 0.695$ et $24305/32768 \approx 0.742$. À chaque étape, le seuil est la valeur du jeu vers lequel on passerait~: $\tfrac12$, puis $\tfrac58$, puis $\tfrac{89}{128}$.

Si le jeu vaut $v$ après un rejet, on garde tout lancer supérieur à $v$. La nouvelle valeur est donc
\[
E[\max(U, v)] = v \cdot v + \int_v^1 x\,dx = \frac{1 + v^2}{2}.
\]
On part de $v = \tfrac12$, la valeur d'un dernier tirage qu'il faut obligatoirement garder.

La récurrence admet pour seul point fixe $v = 1$, donc avec beaucoup de tirages la valeur tend vers le haut de l'intervalle, comme il se doit.

\iqwrong{Donner $2/3$ pour un rejet, ce qui est la moyenne du plus grand de deux tirages. C'est ce qu'on obtient avec le recul. Ici, il faut décider avant de voir le second tirage.}
\end{solution}

\enlargethispage{-3\baselineskip}%
\begin{solution}{iq:iv1:optimal-stopping:3}
On garde 51 ou plus, et le jeu vaut environ $\$63.13$.

Si l'on détient un ticket de numéro $x$, le second tirage est uniforme sur les 99 autres, donc sa moyenne vaut $(5050 - x)/99$. On garde lorsque
\[
x \ge \frac{5050 - x}{99} \iff 100x \ge 5050 \iff x \ge 50.5.
\]
Le seuil est exactement la moyenne simple, $50.5$. Retirer un ticket bas relève la moyenne des autres, en retirer un haut la fait baisser, et les deux effets se compensent à cet endroit.

La valeur est $\frac{1}{100}\bigl(\sum_{x = 51}^{100} x + \sum_{x = 1}^{50} \frac{5050 - x}{99}\bigr) = \frac{1}{100}(3775 + 2537.6)$. La version continue, avec un seuil de $\tfrac12$, vaut $0.625$, ce qui est proche de ce résultat à l'échelle de 100.

\iqwrong{Croire qu'il faut changer seulement pour les tickets très en dessous du milieu, à cause de «~le risque de perdre ce que l'on a~». On reçoit une moyenne, et un joueur neutre au risque doit utiliser le seuil égal à la moyenne.}
\end{solution}

\begin{solution}{iq:iv1:optimal-stopping:4}
Aucune règle ne fait mieux que $1/2$, et toute règle atteint exactement $1/2$.

La proportion de rouges parmi les cartes restantes est une martingale\index{martingale}. Si $r$ des $r + b$ cartes restantes sont rouges, la proportion espérée après la carte suivante est
\[
\frac{r}{r + b} \cdot \frac{r - 1}{r + b - 1} + \frac{b}{r + b} \cdot \frac{r}{r + b - 1} = \frac{r}{r + b}.
\]
La chance de gagner lorsqu'on s'arrête est la proportion courante. Le théorème d'arrêt facultatif\index{arrêt facultatif} affirme que l'espérance à toute règle d'arrêt vaut la valeur de départ, $\tfrac12$.

Une preuve plus courte~: sachant ce qu'on a vu, la carte suivante et la toute dernière carte sont échangeables. Arrêter et parier sur la carte suivante a donc la même chance que parier sur la dernière carte du jeu, et celle-ci est rouge avec une chance de $\tfrac12$, quelle que soit la règle.

\iqwrong{Attendre que davantage de cartes noires soient passées, car le jeu est alors plus riche en rouges. Il s'enrichit effectivement quand cela arrive, mais cela peut ne jamais arriver, et en moyenne ces trajectoires se compensent. Une induction à rebours montre que chaque état est exactement indifférent.}
\iqwhy{C'est l'exemple le plus simple de «~aucune règle de choix du moment ne bat un jeu équitable~». Quiconque présente une règle qui semble tirer un avantage du choix du moment sur un processus équitable a commis une erreur, ou sélectionne son échantillon.}
\end{solution}

\begin{solution}{iq:iv1:optimal-stopping:5}
Pour $n = 3$~: rejeter le premier, puis prendre le premier candidat meilleur que tous les précédents, avec une chance de gagner de $1/2$. Pour $n = 10$~: rejeter les 3 premiers, avec une chance de gagner d'environ $0.399$. Pour $n$ grand~: rejeter une fraction $1/e$, avec une chance de gagner qui tend vers $1/e \approx 0.368$.

Une bonne règle rejette les $s$ premiers, puis prend le candidat suivant qui est meilleur que tous les précédents\index{problème du secrétaire}. On gagne si le meilleur global est à la position $k > s$, avec une chance de $\tfrac1n$ pour chaque $k$, et si le meilleur des $k - 1$ premiers se trouve parmi les $s$ premiers, avec une chance de $\tfrac{s}{k - 1}$. Donc
\[
P(\text{gain}) = \frac{s}{n} \sum_{k = s + 1}^{n} \frac{1}{k - 1}.
\]
Pour $n = 3$ et $s = 1$, cela vaut $\tfrac13(1 + \tfrac12) = \tfrac12$. Prendre n'importe quel candidat à position fixe ne gagne qu'avec une chance de $\tfrac13$. Pour $n$ grand, avec $s = xn$, la somme vaut environ $\ln(1/x)$, donc on maximise $x\ln(1/x)$, dont le maximum est atteint en $x = 1/e$ et vaut $1/e$.

Il n'est pas nécessaire de supposer la forme à seuil. À la position $k$, un candidat qui bat tous les précédents est le meilleur global avec une chance de $k/n$, chance qui croît avec $k$, tandis que la valeur de la poursuite diminue. Une induction à rebours sur $k$ donne la même réponse, $0.399$ pour $n = 10$.

\iqwrong{Supposer que la chance de gagner tend vers 0 quand $n$ croît, parce que le meilleur est de plus en plus difficile à trouver. La chance de gagner tend vers une constante, quel que soit le nombre de candidats.}
\end{solution}

\begin{solution}{iq:iv1:optimal-stopping:6}
On s'arrête au premier nombre supérieur à $0.8$. On effectue 5 tirages en moyenne, et le jeu vaut $0.8$. En général, le seuil vaut $1 - \sqrt{2c}$.

Soit $v$ la valeur du jeu avant un tirage qu'il faut payer\index{prix de réservation}. L'avenir paraît identique après tout rejet, donc un tirage $x$ est gardé lorsque $x \ge v$. L'équation de la valeur est $v = E[\max(U, v)] - c$, qui s'écrit
\[
\int_v^1 (x - v)\,dx = c, \quad\text{donc}\quad \frac{(1 - v)^2}{2} = c.
\]
Avec $c = 0.02$, on obtient $1 - v = 0.2$. Le nombre de tirages suit une loi géométrique de probabilité de succès $0.2$, donc sa moyenne vaut 5. Par un second calcul~: on obtient un nombre de moyenne $0.9$, et l'on paie $5 \cdot 0.02$, ce qui laisse $0.8$.

Le seuil baisse quand le coût augmente. Pour $c = \tfrac12$, le seuil vaut 0 et le jeu ne vaut rien, car le premier tirage ne vaut que son coût. Pour $c = 0$, le seuil vaut 1, et l'on poursuit indéfiniment.

\iqwrong{Abaisser le seuil à mesure que les coûts s'accumulent. L'argent déjà dépensé est perdu, et l'avenir paraît identique après chaque tirage, donc il n'y a qu'un seul seuil pour tout le jeu.}
\iqwhy{Continuer est une option qui a un prix. Cette même équation apparaît lorsqu'on décide combien de temps chercher un meilleur prix, ou quand arrêter de collecter des données.}
\end{solution}

\begin{solution}{iq:iv1:optimal-stopping:7}
On garde 15 ou plus. Le jeu vaut $91/6 \approx \$15.17$.

Soit $w$ la valeur du jeu lorsqu'on s'apprête à faire un lancer gratuit. Une relance rapporte donc $w - 1$, si bien qu'on garde tout nombre supérieur à $w - 1$. Une règle à seuil «~garder au moins $t$~» donne un nombre gardé de moyenne $(t + 20)/2$, et $20/(21 - t)$ lancers en moyenne, dont tous sauf le premier coûtent 1~\$. Sa valeur vaut
\[
\frac{t + 20}{2} - \Bigl(\frac{20}{21 - t} - 1\Bigr).
\]
On obtient $15.14$ pour $t = 14$, $15.17$ pour $t = 15$, et $15$ pour $t = 16$. Donc $t = 15$ est le meilleur. Cela concorde avec la condition~: $w - 1 = 14.17$, donc on garde ce qui dépasse $14.17$, soit 15 ou plus. L'itération sur la valeur, $w = E[\max(k, w - 1)]$, atteint la même valeur.

Cas limites~: avec une relance gratuite, la valeur serait de 20. Avec un coût très élevé, on garde le premier lancer, pour une valeur de $10.5$.

\iqwrong{Comparer le lancer avec $w$ lui-même, ce qui donne un seuil de 16 et une valeur de seulement 15. Une relance coûte 1~\$, donc le nombre à battre est $w - 1$.}
\end{solution}

\begin{solution}{iq:iv1:optimal-stopping:8}
On accepte toute offre égale ou supérieure à $0.5$. L'attente dure 2 jours en moyenne, et le bloc vaut $0.625$ juste avant la première offre.

Soit $v$ la valeur avant l'arrivée d'une offre\index{recherche actualisée}. Une offre refusée vaut $\delta v$, avec $\delta = 0.8$, donc on accepte lorsque l'offre dépasse $\delta v$. Alors
\[
v = E[\max(U, \delta v)] = \frac{1 + (\delta v)^2}{2}.
\]
En résolvant cette équation, on obtient $v = \frac{1 - \sqrt{1 - \delta^2}}{\delta^2} = \frac{1 - 0.6}{0.64} = 0.625$, et le seuil vaut $\delta v = 0.5$. La moitié des offres sont acceptées, donc l'attente suit une loi géométrique, de moyenne 2 jours. Une simulation avec cette règle donne la même valeur.

Les cas limites sont cohérents. Pour $\delta = 0$, on accepte toute offre, ce qui vaut $\tfrac12$. Pour $\delta = 1$, on attend une offre proche de 1, et la valeur tend vers 1.

\iqwrong{Attendre un seuil de $0.625$, la valeur du bloc, règle qui ne vaut qu'environ $0.609$. Ce qu'on abandonne en refusant, c'est cette valeur un jour plus tard, donc le seuil est $\delta v$, et non $v$. Le fait que le seuil soit égal à l'offre moyenne, $0.5$, est une coïncidence liée à $\delta = 0.8$.}
\end{solution}

\begin{solution}{iq:iv1:optimal-stopping:9}
On relance tant que le total est inférieur à 15, et l'on s'arrête à 15 ou plus (on est indifférent à 15 exactement). Le jeu vaut environ $\$6.15$.

Avec un total de $S$, un lancer de plus a pour espérance~:
\[
\frac16 \sum_{k = 1}^{5} (S + k) = \frac{5S}{6} + \frac{15}{6}.
\]
Cela l'emporte sur $S$ exactement lorsque $S < 15$. Une règle qui compare l'arrêt à un pas de plus, puis l'arrêt, est souvent appelée anticipation à un pas\index{anticipation à un pas}. Elle est optimale ici, car, dès qu'elle dit «~arrêter~», elle le dit encore à chaque total plus élevé, et les totaux ne font que croître.

Pour la valeur, on procède par récurrence rétrograde à partir des totaux de 15 et plus, où la valeur est égale au total. Pour $S < 15$, la valeur vaut $\frac16\sum_{k = 1}^5 V(S + k)$. En descendant, on obtient $V(0) \approx 6.15$. Une simulation de la règle concorde.

On voit combien coûte le risque d'obtenir un zéro. Même si l'on suit la règle, on n'atteint pas 15 dans environ 62~\% des cas, et toute la valeur provient des trajectoires qui y parviennent.

\iqwrong{S'arrêter lorsque le total atteint $3.5 \cdot 5 = 17.5$, ou après un nombre fixé de lancers. La règle dépend de ce que l'on détient, car ce que l'on risque est le total, alors que le gain moyen est fixe, égal à $2.5$, de sorte que le risque et le gain s'équilibrent à 15.}
\end{solution}

\begin{solution}{iq:iv1:optimal-stopping:10}
Si les gains négatifs sont autorisés, toute règle vaut 0. Avec le plancher à 0, aller jusqu'au bout est optimal~: le jeu vaut $3/4$.

La différence pile moins face, $S_k$, est une marche aléatoire équilibrée, donc une martingale\index{arrêt optionnel}. Pour un temps d'arrêt borné, $E[S_\tau] = S_0 = 0$. Toutes les règles valent donc la même chose~: s'arrêter en avance rapporte des gains, compensés par les pertes des trajectoires qui ont dû aller jusqu'au bout.

Avec un plancher, le gain vaut $\max(S_k, 0)$, fonction convexe d'une martingale, donc une sous-martingale\index{sous-martingale}. Son espérance ne peut que croître avec le temps, donc attendre la fin est optimal. Alors
\[
E[S_4^+] = \frac{1 \cdot 4 + 4 \cdot 2}{16} = \frac34,
\]
en distinguant $S_4 = 4$, de probabilité $\tfrac{1}{16}$, et $S_4 = 2$, de probabilité $\tfrac{4}{16}$. Pour $n$ grand, ceci croît comme $\sqrt{n/2\pi}$.

\iqwrong{S'arrêter sur un gain de 1 au premier lancer pour verrouiller un profit, le plancher étant en place. S'arrêter rapporte 1, et continuer rapporte $5/4$, car le plancher protège de la chute. (S'arrêter à 2 après deux piles ne coûte rien par hasard, puisque la marche ne peut plus descendre sous 0.)}
\iqwhy{C'est pourquoi il n'est pas optimal d'exercer par anticipation une option d'achat américaine sur une action sans dividende et à taux nuls. Le gain de l'option d'achat est convexe dans une martingale, donc l'exercice anticipé revient à abandonner la valeur restante de l'option.}
\end{solution}

\begin{solution}{iq:iv1:optimal-stopping:11}
Le rapport vaut au plus 10. Un rapport parfait vaut au plus 20.

Sans rapport\index{valeur de l'information}, on négocie, avec une espérance $0.5 \cdot 100 - 40 = 10$.

Avec le rapport, «~oui~» et «~non~» sont équiprobables, par symétrie. Après «~oui~», la probabilité d'un accord est $0.8$, donc négocier vaut $40$. Après «~non~», elle vaut $0.2$, donc négocier vaut $-20$, et l'on s'abstient, ce qui donne 0. La valeur est donc $\tfrac12 \cdot 40 + \tfrac12 \cdot 0 = 20$, soit 10 de plus que la valeur initiale.

Avec un rapport parfait, on ne négocie que lorsque l'accord va se conclure, et l'on gagne 60. La valeur vaut donc $\tfrac12 \cdot 60 = 30$, soit 20 de plus.

La valeur ne vient que du changement de décision. Un rapport juste dans 60~\% des cas ne vaut rien ici~: après «~non~», on aurait une probabilité de $0.4$, qui vaut 0, donc il ne permet d'éviter aucune perte. C'est un contrôle important~: une précision supérieure à $\tfrac12$ ne signifie pas que l'information a de la valeur.

\iqwrong{Retrancher le prix au profit espéré plus élevé après «~oui~», soit $40 - 10 = 30$. On n'obtient pas «~oui~» à chaque fois, et il faut s'engager sur le prix avant de savoir ce que dit le rapport.}
\end{solution}

\begin{solution}{iq:iv1:optimal-stopping:12}
Lancer d'abord la pièce B, puis rester avec B après pile et passer à A après face. Le jeu vaut $13/12$, contre 1 pour deux lancers de la pièce A.

La pièce B a la même probabilité au premier lancer, $\tfrac12$. Mais elle apporte aussi une information\index{exploration ou exploitation}. Après pile, la probabilité pour B est $\tfrac23$ selon la règle de Laplace~; après face, elle est $\tfrac13$, donc on passe à A. La valeur est donc
\[
\tfrac12 + \tfrac12 \cdot \tfrac23 + \tfrac12 \cdot \tfrac12 = \tfrac{13}{12}.
\]
Le gain supplémentaire de $1/12$ est la valeur de l'apprentissage. Avec un seul lancer restant, il n'y a rien à apprendre, et les deux pièces valent $\tfrac12$. Avec trois lancers, la valeur est $5/3$, contre $3/2$ pour la pièce A, soit un edge de $1/6$. La valeur de l'exploration croît avec le temps qui reste pour exploiter ce que l'on apprend.

\iqwrong{Dire qu'on est indifférent parce que les deux pièces ont la même moyenne. La pièce inconnue a la même moyenne, mais elle a une variance sur ce qu'on saura ensuite, et l'on peut l'exploiter.}
\iqwhy{Pour deux stratégies de même rendement espéré, il faut affecter le premier capital à celle que l'on connaît le moins, tant qu'on disposera du temps nécessaire pour agir sur ce qu'on apprend.}
\end{solution}

\begin{solution}{iq:iv1:optimal-stopping:13}
On s'arrête dès que le total atteint au moins $\sqrt{2} - 1 \approx 0.414$. Le jeu vaut $(\sqrt{2} - 1)e^{\sqrt{2} - 1} \approx 0.627$ au départ.

Avec un total de $s$, un tirage de plus puis l'arrêt donnent une espérance de
\[
\int_0^{1 - s} (s + u)\,du = s(1 - s) + \frac{(1 - s)^2}{2} = \frac{1 - s^2}{2}.
\]
Cela l'emporte sur $s$ exactement lorsque $s^2 + 2s - 1 < 0$, donc lorsque $s < \sqrt{2} - 1$\index{anticipation à un pas}. Une fois ce seuil dépassé, le total ne fait que croître, donc la règle reste «~arrêter~».

Pour la valeur en $s$ sous le seuil $t$, la fonction valeur vérifie $V(s) = \int_s^1 V(x)\,dx$, où $V(x) = x$ au-dessus de $t$. En dérivant, on obtient $V' = -V$, donc $V(s) = V(t)e^{t - s}$, et $V(t) = t$. Donc $V(0) = te^t$. Une simulation de la règle donne la même valeur, et un balayage d'autres seuils confirme le maximum en $t$.

\iqwrong{S'arrêter dès que le total dépasse $\tfrac12$, par symétrie. Il n'y a pas de symétrie, car dépasser coûte tout ce que l'on possède, et à $\tfrac12$ cela représente beaucoup pour une faible probabilité d'un gain modeste.}
\end{solution}

\begin{solution}{iq:iv1:optimal-stopping:14}
Le jeu composé de 2 rouges et 2 noires vaut $2/3$, et le jeu complet vaut environ 2.62~\$.

Soit $V(r, b)$ la valeur lorsqu'il reste $r$ rouges et $b$ noires. Alors $V(0, b) = 0$, $V(r, 0) = r$ et\index{recurrence@récurrence rétrograde}
\[
V(r, b) = \max\Bigl(0,\ \frac{r}{r + b}\bigl(1 + V(r - 1, b)\bigr) + \frac{b}{r + b}\bigl(V(r, b - 1) - 1\bigr)\Bigr).
\]
Il s'agit d'un tableau de $27 \times 27$, que l'on remplit à partir du petit coin. On s'arrête lorsque le second terme n'est pas positif.

À la main~: $V(1, 1) = \tfrac12$. $V(1, 2)$ a pour valeur de continuation $\tfrac13(1 + 0) + \tfrac23(\tfrac12 - 1) = 0$. Et $V(2, 1) = \tfrac23(1 + \tfrac12) + \tfrac13(2 - 1) = \tfrac43$. Donc
\[
V(2, 2) = \tfrac12\bigl(1 + V(1, 2)\bigr) + \tfrac12\bigl(V(2, 1) - 1\bigr) = \tfrac12 + \tfrac16 = \tfrac23.
\]
Le jeu complet vaut 0 si l'on ne commence jamais, et il se termine à zéro si l'on ne s'arrête pas, donc toute la valeur tient à la possibilité de s'arrêter\index{valeur de l'option d'arrêt}. On remarque que le jeu a une valeur, alors même que chaque carte est un pari équitable. Un prophète qui s'arrêterait au sommet de la marche gagnerait environ 4 sur un jeu complet, et le jeu en obtient environ les deux tiers sans aucune prescience.

\iqwrong{Dire que la valeur est 0, par le théorème d'arrêt optionnel. Le total courant n'est pas une martingale~: lorsqu'on est en avance, il reste plus de cartes noires que de rouges, donc la carte suivante a une espérance négative. C'est précisément à ce moment-là qu'on s'arrête.}
\end{solution}

\begin{solution}{iq:iv1:optimal-stopping:15}
On prend la première place libre située à au plus 6 unités avant le restaurant, ou à défaut la première place libre à partir de celui-ci. La marche espérée vaut environ $6.57$. En général, avec $q = 1 - p$, on prend une place libre située $x$ unités avant le restaurant exactement lorsque $q^x \ge \tfrac12$.

Utiliser la récurrence rétrograde\index{recurrence@récurrence rétrograde}. Au-delà du restaurant, la marche ne fait que croître, donc on prend la première place libre, située en moyenne à $1/p$ unités au-delà. Soit $V(x)$ la marche espérée en arrivant $x$ unités avant le restaurant, avant de voir cette place. Alors
\[
V(0) = \frac{q}{p}, \qquad V(x) = p\min\bigl(x, V(x - 1)\bigr) + q\,V(x - 1).
\]
Si l'on prend toute place libre à partir de $x$ vers l'intérieur, cela se résout en $V(x) = x + (2q^{x + 1} - q)/p$. À $x$ unités, prendre la place coûte $x$, et passer coûte $V(x - 1) = x - 1 + (2q^x - q)/p$. On prend donc la place lorsque $p \le 2q^x - q$, soit $q^x \ge \tfrac12$\index{anticipation à un pas}. En clair~: on la prend dès qu'il y a au moins une chance sur deux que les $x$ places, de la place courante jusqu'à celle du restaurant incluse, soient toutes occupées.

Avec $q = 0.9$, on a $0.9^6 = 0.531$ et $0.9^7 = 0.478$, donc la limite est de 6 unités, et $V(6) = 6 + (2 \cdot 0.9^7 - 0.9)/0.1 \approx 6.57$. Des limites de 5 et de 7 unités donnent $6.63$ et $6.61$.

\iqwrong{Prendre la première place libre dans un rayon de 10 unités, car l'écart moyen entre places libres vaut $1/p = 10$. Cela coûte $7.28$. Le bon test compare une probabilité à un demi, et non une distance à l'écart moyen. Refuser toute place avant le restaurant coûte $9$.}
\iqwhy{Dès qu'on s'arrêterait en un point, on s'arrêterait aussi en tout point plus proche, donc comparer l'arrêt immédiat avec l'arrêt à la prochaine occasion suffit. Des règles de ce type monotone apparaissent lorsqu'on traite un ordre avant une échéance, et elles évitent de résoudre toute la récurrence.}
\end{solution}

\begin{solution}{iq:iv1:optimal-stopping:16}
Les profits valent 0, $\tfrac13$, $\tfrac13$ et $\tfrac25$. Le lancer gratuit supplémentaire n'ajoute rien, et la limite vaut $\tfrac12$. Pour $n$ impair, la valeur est $\frac{n + 1}{2(n + 2)}$, et chaque $n$ pair a la même valeur que le nombre impair qui le précède.

Avec $k$ piles en $n$ lancers, la probabilité d'avoir pile au lancer suivant est $\frac{k + 1}{n + 2}$, selon la règle de Laplace\index{règle de succession de Laplace}. On mise du côté le plus probable, donc le profit vaut $|2\frac{k + 1}{n + 2} - 1|$. Avec une loi a priori uniforme, $k$ est uniforme sur $0, \dots, n$. La valeur est donc la moyenne de cette valeur absolue sur $k$.

Pour $n = 1$~: probabilités $\tfrac13, \tfrac23$, et profit $\tfrac13$.

Pour $n = 2$~: probabilités $\tfrac14, \tfrac12, \tfrac34$, et profits $\tfrac12, 0, \tfrac12$, donc la moyenne vaut $\tfrac13$.

Pour $n = 3$~: profits $\tfrac35, \tfrac15, \tfrac15, \tfrac35$, donc la moyenne vaut $\tfrac25$.

Miser du côté du premier lancer reste une règle optimale après le deuxième lancer. Si les deux lancers concordent, on mise de ce côté de toute façon, et s'ils divergent, on est indifférent. Le deuxième lancer ne change donc aucune décision et n'apporte aucune valeur. La limite vaut $E|2p - 1| = \tfrac12$, ce que vaudrait la connaissance exacte de $p$.

\iqwrong{Supposer que la valeur de l'information croît régulièrement avec $n$. Elle reste ici constante entre une étape impaire et une étape paire, car seul un changement de décision a de la valeur. Une donnée supplémentaire qui affine l'estimation sans modifier le pari n'ajoute rien à un pari unique de taille fixe.}
\end{solution}

\begin{solution}{iq:iv1:optimal-stopping:17}
A doit garder un premier tirage d'au moins $1/\sqrt{3} \approx 0.577$, et la probabilité pour A de gagner vaut environ $0.462$.

Commençons par B\index{induction arrière}. Étant donné le nombre final $x$ de A, B gagne sauf si ses deux tirages sont tous deux inférieurs à $x$, car B relance exactement lorsqu'il est derrière. A gagne donc avec la probabilité $x^2$.

Considérons maintenant que A dispose d'un tirage $x$. Garder donne $x^2$. Relancer donne un nombre uniforme, et donc $E[U^2] = \tfrac13$. A garde donc son tirage lorsque $x^2 \ge \tfrac13$. Avec un seuil $t$, la probabilité que A gagne vaut
\[
\int_t^1 x^2\,dx + t \cdot \tfrac13 = \frac{1 - t^3}{3} + \frac{t}{3},
\]
qui est maximale pour $t^2 = \tfrac13$, et vaut $\tfrac13\bigl(1 + \tfrac{2}{3\sqrt{3}}\bigr) \approx 0.462$. Pour vérifier, $t = 0$ et $t = 1$ reviennent tous deux à ce que A dispose en pratique d'un seul tirage, ce qui donne $\tfrac13$. Une simulation du jeu donne le même résultat.

\iqwrong{Prendre $\tfrac12$ comme seuil, ce qui maximise le nombre espéré. A ne cherche pas à obtenir un nombre élevé en moyenne~; A doit survivre aux deux tentatives de B. Le gain vaut $x^2$, et non $x$. La convexité favorise un choix plus risqué.}
\end{solution}

\begin{solution}{iq:iv1:optimal-stopping:18}
S'arrêter au premier succès parmi les 9 derniers essais. La probabilité de gagner vaut $0.9^9 \approx 0.387$. L'embauche revient au même problème, avec une probabilité de succès $1/k$ à l'essai $k$.

S'arrêter sur un succès à l'essai $k$ fait gagner si les $n - k$ essais restants sont tous des échecs. Un succès vaut donc d'autant plus qu'il arrive tard, et la règle est un seuil sur le numéro de l'essai. Pour une règle qui accepte tout succès parmi les $s$ derniers essais, on gagne lorsqu'exactement un succès survient parmi eux~:
\[
P(\text{gain}) = s p q^{s - 1},\quad q = 1 - p.
\]
En comparant $s$ et $s + 1$, la probabilité croît tant que $s \cdot p/q < 1$. Ici $p/q = 1/9$, donc $s = 9$ donne exactement 1, et $s = 9$ et $s = 10$ sont à égalité, à $0.387$. Pour $s = 8$, on obtient $0.383$.

C'est l'algorithme des cotes\index{algorithme des cotes}~: on additionne les cotes de succès en remontant depuis la fin, et on commence à accepter lorsque la somme atteint 1. Dans le problème de l'embauche, «~le candidat $k$ bat tous les précédents~» est un succès de probabilité $\frac1k$, ces événements sont indépendants, et le dernier succès est le meilleur de l'ensemble. Les cotes valent $\frac{1}{k - 1}$. Pour $n = 10$, la somme prise depuis la fin atteint 1 en $k = 4$, donc on rejette les 3 premiers, comme le donne le calcul direct. La probabilité de gagner vaut au moins environ $1/e$ lorsque la somme des cotes atteint 1.

\iqwrong{Attendre que le nombre espéré de succès restants soit d'environ 1, soit les 10 derniers essais ici. C'est presque juste, et cela donne la même probabilité de gain dans ce cas d'égalité. En général, il faut utiliser les cotes et non les probabilités, ce qui compte lorsque les probabilités ne sont pas petites.}
\end{solution}

\begin{solution}{iq:iv1:optimal-stopping:19}
La règle «~s'arrêter au premier avantage~» vaut exactement $\pi/4 \approx 0.785$. L'induction arrière sur un long horizon donne environ $0.793$ pour la meilleure règle, donc cette règle n'est pas optimale. Il ne faut jamais s'arrêter lorsqu'on est à égalité ou derrière.

Pour la borne inférieure~: une pile au premier lancer rapporte 1. Après un premier face, la marche équilibrée revient à égalité avec probabilité 1, on peut donc toujours atteindre une fraction de $\tfrac12$. On obtient au moins $\tfrac12 \cdot 1 + \tfrac12 \cdot \tfrac12 = \tfrac34$.

Pour le premier avantage\index{jeu de Chow--Robbins}~: le premier avantage survient au lancer $2m + 1$ avec probabilité $\frac{C_m}{2^{2m + 1}}$, où $C_m$ désigne le nombre de Catalan, et la fraction vaut alors $\frac{m + 1}{2m + 1}$. La valeur est donc
\[
\sum_{m \ge 0} \frac{C_m}{2^{2m + 1}} \cdot \frac{m + 1}{2m + 1} = \sum_{m \ge 0} \binom{2m}{m}\frac{1}{(2m + 1)2^{2m + 1}} = \frac12 \arcsin 1 = \frac{\pi}{4},
\]
en utilisant $C_m(m + 1) = \binom{2m}{m}$ et la série de l'arcsinus\index{série de l'arcsinus}.

Il ne faut jamais s'arrêter lorsqu'on est à égalité ou derrière. Depuis n'importe quel état, on peut revenir à au moins $\tfrac12$ presque sûrement, car la marche est récurrente, donc s'arrêter en dessous de $\tfrac12$ fait perdre. Exactement à $\tfrac12$, on a une probabilité positive de prendre un avantage, et on peut toujours revenir, donc continuer est strictement meilleur.

Le premier avantage n'est pas toujours un bon moment pour s'arrêter non plus. Après pile, face, face, on a $\tfrac23$, mais continuer vaut environ $0.669$. La règle optimale exige donc un avantage plus grand à certains moments, et l'avantage requis croît comme $0.84\sqrt{n}$. La valeur provient d'une induction arrière depuis 5000 lancers, avec une estimation basse et une estimation haute de la valeur terminale, et les deux arrondissent à $0.793$. C'est le résultat de ce calcul, non d'une formule.

\iqwrong{Affirmer que le jeu vaut 1 parce qu'une marche équilibrée finit par prendre une avance quelconque. Elle prend bien de l'avance, mais c'est la fraction qui rapporte, et à mesure que le nombre de lancers croît, une avance donnée vaut moins. Être en avance d'un montant fixe à un instant tardif a une valeur proche de $\tfrac12$.}
\end{solution}

\begin{solution}{iq:iv1:optimal-stopping:20}
Prenons le seuil $t = \tfrac12 E[M]$, avec $M = \max X_i$, et acceptons le premier prix au moins égal à $t$. On gagne alors au moins $\tfrac12 E[M]$. Un exemple avec un prix de 1 certain, suivi d'un prix de $1/\varepsilon$ avec probabilité $\varepsilon$, montre qu'on ne peut pas faire mieux que la moitié en général.

Soit $p = P(M \ge t)$, la probabilité d'accepter quelque chose\index{inégalité du prophète}. Le gain comporte deux parties. Chaque fois qu'on accepte, on obtient au moins $t$, ce qui donne $pt$. À cela s'ajoute l'excédent $(X_i - t)^+$ lorsqu'on accepte le prix $i$. On atteint le prix $i$ sans avoir encore accepté avec une probabilité d'au moins $1 - p$, et cet événement ne dépend que des prix précédant $i$, il est donc indépendant de $X_i$. Ainsi
\[
\text{Gain} \ge pt + (1 - p)\sum_i E[(X_i - t)^+] \ge pt + (1 - p) E[(M - t)^+] \ge pt + (1 - p)(E[M] - t),
\]
où la somme est supérieure à $(M - t)^+$ car le plus grand excédent figure parmi les termes. Avec $t = \tfrac12 E[M]$, $t$ et $E[M] - t$ valent tous deux $\tfrac12 E[M]$, et $p$ disparaît~:
\[
\text{Gain} \ge p \cdot \tfrac12 E[M] + (1 - p) \cdot \tfrac12 E[M] = \tfrac12 E[M].
\]
Pour l'exemple, le premier prix vaut 1 avec certitude. Le second prix vaut $1/\varepsilon$ avec probabilité $\varepsilon$, donc son espérance vaut elle aussi 1. Toute stratégie vaut donc 1. Mais le prophète prend le second prix lorsqu'il apparaît, et sinon le premier, pour $2 - \varepsilon$. Le rapport tend vers $\tfrac12$.

\iqwrong{Chercher un seuil dépendant de manière compliquée de l'ensemble des distributions. Un seul nombre donne déjà la meilleure garantie possible sur toutes les distributions, et la preuve n'utilise rien de l'ordre dans lequel elles arrivent.}
\iqwhy{Cela donne la véritable valeur de l'information a posteriori dans les problèmes de sélection. Les backtests qui choisissent après coup le meilleur moment peuvent afficher jusqu'au double de ce que toute règle réelle peut rapporter, même lorsque les distributions sont parfaitement connues.}
\end{solution}
